% ============================================================
% Chapter 2 — System Requirements Specification (SRS)
% Source: PSYCON Engineering Design Document, Vol. 1, Ch. 2
% ============================================================
\section{System Requirements Specification}
\label{sec:srs}

\subsection{Purpose}
\label{sec:srs-purpose}

This section defines exactly what PSYCON is expected to do from an engineering perspective. The requirements stated here act as the baseline against which the prototype will be evaluated.

\subsection{Design Objectives}
\label{sec:design-objectives}

The system shall:
\begin{itemize}
    \item Continuously collect physiological data.
    \item Continuously collect speech samples.
    \item Continuously monitor environmental conditions.
    \item Synchronize all collected data.
    \item Operate as a wearable prototype.
    \item Be modular.
    \item Be expandable.
    \item Be low-cost.
    \item Be reproducible using commercially available components.
    \item Support future AI-based analysis.
\end{itemize}

\subsection{Functional Requirements}
\label{sec:functional-requirements}

\subsubsection{FR-1: Physiological Monitoring}
The wrist module shall measure heart rate, pulse waveform (PPG), skin conductance (EDA/GSR), skin temperature, wrist motion, and wrist orientation.

\noindent\textit{Status:} \doneTag\ Hardware selected \quad \pendingTag\ Integration pending

\subsubsection{FR-2: Audio Monitoring}
The audio module shall capture speech, digitize audio using I\textsuperscript{2}S, stream or store recordings, and maintain synchronized timestamps.

\noindent\textit{Status:} \doneTag\ Hardware selected \quad \pendingTag\ Firmware pending

\subsubsection{FR-3: Environmental Monitoring}
The system shall measure ambient light for purposes of context awareness, sleep/environment correlation, and future feature extraction.

\noindent\textit{Status:} \doneTag\ Hardware selected

\subsubsection{FR-4: Data Synchronization}
The system shall timestamp every sample, preserve acquisition order, and allow synchronization between the wrist and audio modules. This requirement is critical because machine-learning analysis becomes unreliable if physiological and audio events are misaligned in time.

\noindent\textit{Status:} \pendingTag\ Firmware requirement

\subsubsection{FR-5: Wireless Communication}
Each ESP32 shall support Wi-Fi and BLE (future expansion), enabling data upload, remote monitoring, over-the-air (OTA) firmware updates (future), and local communication.

\noindent\textit{Status:} \doneTag\ Supported by hardware

\subsubsection{FR-6: Local Processing}
Each module shall read sensors, filter signals, package data, detect sensor failures, and continue operation independently of the other module.

\subsubsection{FR-7: Modular Operation}
If one module fails, the second module shall continue operating. For example, an audio-module failure shall not stop physiological acquisition, and a physiological-module failure shall not stop audio acquisition.

\subsubsection{FR-8: Expandability}
Future sensors should be addable without redesigning the architecture. Candidate future sensors include ECG, EEG, EMG, respiration, skin hydration, and air-quality sensing.

\subsection{Non-Functional Requirements}
\label{sec:non-functional-requirements}

\textbf{Reliability.} The target is continuous operation without crashes. The prototype goal is a minimum continuous runtime of 6\,h, with a stretch goal of 24\,h.

\textbf{Accuracy.} Sensors should operate within manufacturer specifications; no artificial software correction should replace proper calibration.

\textbf{Repeatability.} Repeated measurements under identical conditions should produce similar results, which is especially important for GSR, heart rate, and temperature.

\textbf{Maintainability.} The system should use standard breakout boards, use labeled connectors, allow module replacement, and support firmware updates.

\textbf{Portability.} The system should fit on the body, operate on batteries, and avoid wired external power during use.

\textbf{Scalability.} Future revisions should migrate easily to a custom PCB, smaller ESP32 variants, and improved analog front ends.

\textbf{Cost.} The primary objective is high capability with affordable hardware; the approximate prototype cost target is lower than research-grade commercial systems.

\textbf{Safety.} The wearable shall use protected LiPo batteries, avoid exposed battery terminals, maintain a safe GSR excitation current, avoid charging while connected to GSR electrodes, and use insulated wiring.

\subsection{Performance Targets}
\label{sec:performance-targets}

\Cref{tab:performance-targets} summarizes the target operating behavior for each monitored parameter.

\begin{table}[H]
\centering
\caption{Performance targets by parameter.}
\label{tab:performance-targets}
\small
\begin{tabularx}{\columnwidth}{@{}Y l@{}}
\toprule
\textbf{Parameter} & \textbf{Target} \\
\midrule
Heart-rate monitoring & Continuous \\
GSR monitoring & Continuous \\
Temperature monitoring & Continuous \\
Motion monitoring & Continuous \\
Ambient light & Continuous \\
Speech recording & Continuous or scheduled \\
Wireless communication & Stable \\
Battery-powered & Yes \\
Rechargeable & Yes \\
\bottomrule
\end{tabularx}
\end{table}

\subsection{Hardware Constraints}
\label{sec:hardware-constraints}

\Cref{tab:hardware-constraints} lists the hardware constraints of the finalized design, organized by controllers, power, and sensor allocation between the wrist and audio modules.

\begin{table}[H]
\centering
\caption{Hardware constraints of the finalized PSYCON design.}
\label{tab:hardware-constraints}
\small
\begin{tabularx}{\columnwidth}{@{}l Y@{}}
\toprule
\textbf{Category} & \textbf{Components} \\
\midrule
Controllers & 2 $\times$ ESP32 DevKit (30-pin) \\
Power & 2 $\times$ TP4056 USB-C charger; 3 $\times$ 500\,mAh LiPo battery \\
Wrist sensors & MAX30102; ADS1115; MCP9808; MPU6050; GSR electrodes \\
Audio sensors & INMP441; TEMT6000 \\
\bottomrule
\end{tabularx}
\end{table}

\subsection{Software Constraints}
\label{sec:software-constraints}

Firmware should:
\begin{itemize}
    \item Support FreeRTOS tasks (recommended).
    \item Prevent blocking sensor reads.
    \item Handle communication failures gracefully.
    \item Log errors.
    \item Support firmware expansion.
\end{itemize}

\subsection{Data Requirements}
\label{sec:data-requirements}

Every record should include a timestamp, module ID, sensor ID, sensor values, battery level (future), and error flags (future). \Cref{fig:data-record-flow} illustrates the logical structure of a single data record as it is composed prior to storage or transmission.

\begin{figure}[H]
\centering
\begin{tikzpicture}[
    font=\scriptsize,
    node distance=0mm,
    field/.style={draw, rounded corners=2pt, fill=blue!6, minimum width=2.5cm, minimum height=0.65cm, align=center, line width=0.5pt}
]
\node[field] (ts) {Timestamp};
\node[field, below=6mm of ts] (mod) {Module};
\node[field, below=6mm of mod] (sen) {Sensor};
\node[field, below=6mm of sen] (meas) {Measurement};
\node[field, below=6mm of meas] (qf) {Quality Flag};

\draw[-{Stealth[length=2mm]}, line width=0.6pt] (ts.south) -- (mod.north);
\draw[-{Stealth[length=2mm]}, line width=0.6pt] (mod.south) -- (sen.north);
\draw[-{Stealth[length=2mm]}, line width=0.6pt] (sen.south) -- (meas.north);
\draw[-{Stealth[length=2mm]}, line width=0.6pt] (meas.south) -- (qf.north);
\end{tikzpicture}
\caption{Logical structure of a single PSYCON data record, composed sequentially from acquisition timestamp through to a quality flag.}
\label{fig:data-record-flow}
\end{figure}

\subsection{Assumptions}
\label{sec:srs-assumptions}

\begin{assumptionbox}
The following are treated as current engineering assumptions until experimentally verified:
\begin{itemize}
    \item Standard module revisions.
    \item Standard operating voltages.
    \item No I\textsuperscript{2}C address conflicts.
    \item Proper power distribution.
    \item Common ground maintained.
    \item Standard ESP32 DevKit V1 behavior.
\end{itemize}
\end{assumptionbox}

\subsection{Risks}
\label{sec:risks}

\Cref{tab:risks} summarizes known engineering risks by priority level.

\begin{table}[H]
\centering
\caption{Known engineering risks by priority.}
\label{tab:risks}
\small
\begin{tabularx}{\columnwidth}{@{}l Y@{}}
\toprule
\textbf{Priority} & \textbf{Risk} \\
\midrule
\highriskTag & GSR analog front-end quality \\
\highriskTag & Battery runtime \\
\highriskTag & Charging while operating \\
\highriskTag & Long-term stability \\
\midrule
\medriskTag & Audio noise \\
\medriskTag & Motion artifacts \\
\medriskTag & Synchronization drift \\
\midrule
\lowriskTag & I\textsuperscript{2}C conflicts (none expected) \\
\lowriskTag & GPIO availability \\
\lowriskTag & Sensor compatibility \\
\bottomrule
\end{tabularx}
\end{table}

\subsection{Acceptance Criteria}
\label{sec:acceptance-criteria}

The prototype is considered functionally complete when:
\begin{itemize}
    \item All sensors initialize successfully.
    \item No I\textsuperscript{2}C address conflicts occur.
    \item Audio records correctly.
    \item Physiological data streams continuously.
    \item Modules operate independently.
    \item Battery operation is stable.
    \item Data timestamps remain synchronized.
    \item Documentation is complete.
    \item Safety procedures are documented.
    \item Hardware passes integration testing.
\end{itemize}

\subsection{Chapter Status}
\label{sec:ch2-status}

\begin{validationbox}
\textbf{Completed for this chapter:}
\begin{itemize}
    \item Functional requirements
    \item Non-functional requirements
    \item Design objectives
    \item Constraints (hardware and software)
    \item Acceptance criteria
    \item Risk overview
    \item Data requirements
\end{itemize}
\end{validationbox}
